\documentclass[11pt,a4paper]{article}
\usepackage[T1]{fontenc}
\usepackage[utf8]{inputenc}
\usepackage{amsmath,amssymb,amsthm}
\usepackage[margin=1in]{geometry}
\usepackage[colorlinks=true,linkcolor=blue,urlcolor=blue]{hyperref}
\theoremstyle{plain}
\newtheorem{theorem}{Theorem}
\newtheorem{lemma}[theorem]{Lemma}
\newtheorem{proposition}[theorem]{Proposition}
\theoremstyle{definition}
\newtheorem{remark}[theorem]{Remark}

\title{Settling empirical conjectures in the OEIS at scale:\\
\large method, verification, and an honest account of the errors}
\author{Adrian Perez Fontelles\\ \small Independent researcher}
\date{6 September 2026}
\begin{document}
\maketitle

\begin{abstract}
This note describes a body of work that settles 8887 conjectures across 8860 entries of the
On-Line Encyclopedia of Integer Sequences: 8881 proofs and 6 disproofs, each recorded on its
entry as an open conjecture, an empirical observation, or an unverified formula at the time it
was settled. The mathematical core is not new. Most of these entries count arrays of a fixed
width over a finite alphabet, so the count is a walk count in a finite digraph, the sequence is
$C$-finite, and the question ``does this empirical recurrence hold?'' is DECIDABLE rather than
empirical. What is new is only that the question was actually asked, of every such entry, with
the reading of each entry's English pinned against its own published data before anything was
computed. This note explains the method, the verification discipline, and --- at length,
because it is the part a reader should weigh --- the errors the discipline caught.
\end{abstract}

\noindent\small 2020 Mathematics Subject Classification. 05A15, 11B37, 68Q45, 68W30.\normalsize

\section{What is claimed, and what is not}

The OEIS contains many thousands of entries that carry a line beginning \texttt{Empirical:} or
\texttt{Conjecture:}, typically a linear recurrence fitted to the terms the entry publishes and
never proved. A large family of these, contributed chiefly by R.~H.~Hardin, count arrays of a
fixed width over a fixed finite alphabet subject to a local condition.

The claim of this work is narrow and should be read narrowly.

\begin{itemize}
\item For each of 8860 entries, at least one conjecture recorded as open has been settled, and
a self-contained paper states it, proves it, and reports the range verified.
\item The arguments are, in the great majority of cases, one argument: an array count of this
kind is a walk count in a finite digraph, hence $C$-finite, and whether a proposed
constant-coefficient recurrence annihilates it is a finite exact computation.
\item That argument is classical. It is in Stanley~\cite{stanley}, and any competent reader of
these entries could have applied it.
\end{itemize}

What is \emph{not} claimed: that any individual result is deep, that the method is novel, or
that anything here has been peer reviewed. Where a result leans on a classical theorem --- most
visibly the 41 papers that lean on MacMahon's box formula~\cite{macmahon} --- the paper says so
in its abstract and cites it.

The interesting content, such as it is, lies in two places: that the questions turn out to be
decidable at all, and that the failure mode of a program like this is not the mathematics but
the READING. Almost every error found in the course of this work was a misreading of English,
not a miscalculation.

\section{Why these conjectures are decidable}

Fix a width $W$ and an alphabet $\{0,\dots,m\}$, and consider arrays of $W$ columns and $n$
rows subject to a condition that, for each cell, involves only cells within a bounded number of
rows of it. Say the condition reaches $U$ rows up and $D$ rows down. Then a window of $U+D+1$
consecutive rows decides whether the middle row passes, so the last $U+D$ rows are a sufficient
state.

Take those states as the vertices of a digraph $G$, with an edge from one state to the next
exactly when appending the new row keeps every completed cell legal. An array of $n$ rows is a
walk of $n-c$ steps for a constant $c$ read off the entry's offset, and every array arises from
exactly one walk. Writing $M$ for the adjacency matrix, $\iota$ for the indicator of states
that may begin an array and $\tau$ for those that may end one,
\[
a(n)\;=\;\iota^{\!\top}M^{\,n-c}\tau .
\]
So $a$ satisfies the linear recurrence given by the characteristic polynomial of $M$, and is in
particular $C$-finite.

\begin{lemma}\label{lem:crit}
Let $q(t)=t^{d}-\sum_i c_it^{d-i}$ be the characteristic polynomial of a proposed recurrence,
and put $u_j=\iota^{\!\top}M^{\,j}q(M)\tau$. Then $u_j$ is exactly the residual
$a(j+c+d)-\sum_ic_ia(j+c+d-i)$, and the proposed recurrence holds from some point on if and
only if $u_j=0$ for all large $j$. Moreover $S$ consecutive vanishing $u_j$, where $S$ is the
number of states, force every later one, and the last nonzero $u_j$ pins the threshold exactly.
\end{lemma}

\begin{proof}
The first statement is the identity
$M^{j+d}-\sum_ic_iM^{j+d-i}=M^{j}q(M)$. For the second, Cayley--Hamilton makes $M^{S}$ an
integer combination of $I,M,\dots,M^{S-1}$, so the sequence $(u_j)$ satisfies the monic
recurrence given by the characteristic polynomial of $M$; a run of $S$ zeros therefore forces
all later terms to vanish.
\end{proof}

This is the whole engine. It answers the question, it answers it exactly in integer arithmetic,
and it returns the smallest $n$ beyond which the recurrence holds rather than merely asserting
that it eventually does. The computation is carried out on the vector $q(M)\tau$ rather than on
$M$, so the matrix is never formed.

\begin{remark}
Not every entry is of this shape. Where the array's SHAPE is fixed and the ALPHABET $0..n$ is
what grows, there is no digraph to walk; those counts turn out to be quasi-polynomials, settled
by inclusion--exclusion and then by polynomial division. Where the conjecture is a closed form
rather than a recurrence, the tools are hypergeometric summation and Ore-algebra division. Some
91 distinct arguments appear across the roster; the walk count is simply the commonest.
\end{remark}

\section{The difficulty is the English}

The engine above is unconditional once a model is fixed. Fixing the model means reading a
sentence written by a human being, in a compressed and idiosyncratic style, and turning it into
a predicate. That is where the risk lives, and it is worth being concrete.

\paragraph{``The same population''.}
A224654 counts $(n+2)\times3$ matrices over $\{0,1,2\}$ in which every
$3\times3$ subblock has ``the same population''. Two readings present themselves: the sum of
the nine entries, or the multiset of how often each value occurs. On a BINARY alphabet the two
agree, which is why the sum reading survived every binary entry of the family. Over
$\{0,1,2\}$ it does not: the sum reading gives $102789$ at $n=2$, and the entry publishes
$67797$, which is what the multiset reading gives. Both readings agree at $n=1$, where the
condition is vacuous, so a check on the first term alone would have missed it.

\paragraph{``Its 72 absolute element differences''.}
A234834 asks for $3\times3$ subblocks over $\{0,1,2\}$ whose ``72 absolute element
differences'' sum to $34$. Since $72=9\cdot8$, the phrase names the ORDERED pairs; but the sum
over ordered pairs is even, so $34$ is not achievable and the reading gives $0$. Summing over
the $\binom92=36$ UNORDERED pairs gives $4032$, which is the entry's first term. The entry
counts each pair once and names it twice.

\paragraph{``Diagonally, horizontally or vertically''.}
A189305 counts array permutations in which every cell moves zero or one space ``diagonally,
horizontally or vertically''. Read as the eight king moves plus staying put, the count at $n=2$
is $24$; the entry publishes $14$. Each word names an AXIS --- a pair of opposite offsets ---
and the list names three of them, not eight directions.

In every case the entry's own published terms decided the matter, and decided it immediately.
That is the single most useful fact about this problem: the data is not merely a check on the
arithmetic, it is a check on the reading, and it is a very sharp one. A model built from a
misread sentence reproduces the published terms with probability close to zero.

The discipline that follows is: \textbf{pin the reading against the published data before
writing the engine}, and treat any disagreement as evidence about the reading rather than about
the entry.

\section{Verification}

No result was accepted on one calculation. Four checks were applied.

\begin{enumerate}
\item \textbf{The data gate.} The model must reproduce every term the entry publishes, exactly,
in integer arithmetic, with the index shift read off the entry's offset rather than fitted.
\item \textbf{An independent brute force.} The objects are enumerated a second time directly
from the entry's English --- writing out the arrays and testing the condition cell by cell,
with no transfer matrix and no closed form --- and compared against the published terms. This
is deliberately written against the entry, not against the engine.
\item \textbf{The conjecture on the raw data.} The proposed recurrence is evaluated on the
published terms alone, with no model involved, wherever the proved range and the data overlap.
\item \textbf{The annihilation test.} Lemma~\ref{lem:crit} is carried out over the whole vector
to the Cayley--Hamilton bound, in exact integer arithmetic, not sampled.
\end{enumerate}

Check (2) is the one that matters, and it has a failure mode of its own; see below.

\section{The errors}

A reader assessing work of this size is entitled to the error record, so here it is.

\paragraph{A brute force that shares the engine's misreading is not a check.}
One family of entries writes the array with its width growing, so the natural walk runs across
COLUMNS; but the entries also carry a clause about the order in which values first appear,
stated in ROW major order. An engine that transposes the array must not transpose the clause
with it. One did, and the brute force written alongside it made the identical substitution, so
the two agreed with each other and both matched the entry's published terms. The error was
found only by writing a third count in the entry's own orientation, where four of the eight
entries disagreed --- exactly the four whose neighbour list contains both diagonal words.
\textbf{Independence of a check is a property of how it was written, not of the fact that it
exists.}

\paragraph{A check that disagrees with forty papers at once is the check that is wrong.}
A re-audit reported that forty thresholds were off by exactly one. They were not: the audit's
minimal-order computation counted leading terms the entry never publishes, each of which raises
the minimal order by one. A systematic disagreement of a fixed size, against many results at
once, is a signature of the checker.

\paragraph{Three audits that cried wolf, all with one cause.}
Each was written against a different interface from the one that produced the result: text
extraction from a PDF silently dropped superscripts, so $k^n$ was compared as $kn$; a trailing
comma turned an expression into a one-element tuple; and a check that called an engine's
internal routine skipped the division by a fraction stated in the entry's name, so ``half the
number of \ldots'' came out doubled. \textbf{An audit must go through the same interface the
sweep does, or it is testing a different claim.}

\paragraph{A real error, found by re-audit, in 477 papers.}
An earlier version stated the range over which a recurrence had been proved using a converted
index, and the conversion was wrong. The papers' mathematics was sound and their recurrences
true; the stated ranges were not. All 477 were corrected.

\paragraph{Silent refusals.}
A wrapper called one engine with an argument that engine did not take. The resulting
\texttt{TypeError} was swallowed by a broad exception handler, so every entry of that family
was reported as ``too large to model'' when the model in fact built in seconds. Forty-two
results were sitting behind that one bug. A refusal that looks like a size limit and is
actually a crash is the most expensive kind of error in a program like this, because nothing
about it looks wrong.

\paragraph{Results withdrawn.}
Nine papers were withdrawn as duplicates or as settling something already recorded as settled;
one engine, and its ten papers, were withdrawn on discovering that an existing engine already
performed the identical reduction. Three papers named the wrong contributor, an attribution
hard-coded in a template rather than read from the entry; all were corrected.

\paragraph{A prose error in 66 papers.}
For one family the parser rewrites a negated condition positively, so the negation is already
inside the displayed formula; the surrounding sentence then quoted the entry's own ``no'' and
negated it a second time. The models were unaffected and reproduced every published term, but
sixty-six papers described their object as the opposite of what the entry asks. Corrected and
rebuilt.

\section{Scope, and the walls}

The roster covers 8860 entries by 91 distinct arguments. The largest groups are the plain walk
count (3008 papers), a residual test over a square root (432), $3\times3$ subblock conditions
(301), and counts taken up to relabelling of the alphabet, which are recovered from the labelled
counts by falling-factorial inversion (290).

Several classes resist the method and are recorded as such rather than quietly omitted:

\begin{itemize}
\item $n\times n$ arrays, where both sides grow. There is no fixed width to walk along, and
this is a genuine obstruction rather than a gap in the parser.
\item Entries whose state space is large enough that the annihilation test does not finish, even
after merging states with identical futures --- a reduction that often shrinks a state space by
an order of magnitude and is what makes the wider entries tractable at all.
\item Triangular families whose inclusion--exclusion runs over $2^{|E|}$ subsets, which passes
out of reach at around thirty edges.
\end{itemize}

\section{How to check a single result}

Every paper is self-contained and can be checked independently of everything else here, which
is the point of not letting them refer to one another. To check one:

\begin{enumerate}
\item Read the conjecture as the paper quotes it, and compare it with the live OEIS entry.
\item Read the paper's description of the condition, and satisfy yourself it is what the
entry's name says. This is where an error would be.
\item Enumerate the objects yourself for the first two or three values of $n$ and compare with
the entry's published terms. This is cheap and it tests the modelling.
\item The recurrence check is then a finite exact computation, described in the paper.
\end{enumerate}

Step (3) is the one worth doing. If a paper's model is right, everything after it is
mechanical; if it is wrong, step (3) will say so at once.

\begin{thebibliography}{9}
\bibitem{oeis} The OEIS Foundation, \emph{The On-Line Encyclopedia of Integer Sequences},
\url{https://oeis.org}, 2026.
\bibitem{stanley} R.~P.~Stanley, \emph{Enumerative Combinatorics}, Volume 1, 2nd ed.,
Cambridge University Press, 2011. (Transfer-matrix method, Section 4.7.)
\bibitem{stanley2} R.~P.~Stanley, \emph{Enumerative Combinatorics}, Volume 2, Cambridge
University Press, 1999.
\bibitem{macmahon} P.~A.~MacMahon, \emph{Combinatory Analysis}, Volume 2, Cambridge
University Press, 1916.
\bibitem{flajolet} P.~Flajolet and R.~Sedgewick, \emph{Analytic Combinatorics}, Cambridge
University Press, 2009.
\bibitem{kauers} M.~Kauers and P.~Paule, \emph{The Concrete Tetrahedron}, Springer, 2011.
\bibitem{horn} R.~A.~Horn and C.~R.~Johnson, \emph{Matrix Analysis}, 2nd ed., Cambridge
University Press, 2013.
\end{thebibliography}

\end{document}
